\section{Całka Kurzweila--Henstocka}
\providecommand{\atlasref}[2]{#2}
\label{it:sec:kurzweil-henstock}\label{ma1:ch:kurzweil-henstock}
% Import z AM1-36, zachowane etykiety źródłowe i historyczne IT.
Symbol $\lambda_1$ oznacza miarę Lebesgue'a na prostej. Dla ograniczonej funkcji $F$ i niepustego podprzedziału $I$ piszemy $\osc_I F=\sup_I F-\inf_I F$. Gdy $F$ jest ciągła, a $I$ zwarty, supremum i infimum są osiągane, więc $\osc_I F=\max_I F-\min_I F$. Symbole $\AC$ i $\BV$ oznaczają klasy funkcji bezwzględnie ciągłych i o ograniczonej wariacji; $\Var_I F$ jest supremum sum modułów przyrostów po podziałach $I$. Funkcja $F$ jest bezwzględnie ciągła na $[a,b]$, jeżeli dla każdego $\varepsilon>0$ istnieje $\eta>0$ takie, że dla każdej skończonej rodziny przedziałów $[u_j,v_j]$ o rozłącznych wnętrzach z $\sum_j(v_j-u_j)<\eta$ zachodzi $\sum_j|F(v_j)-F(u_j)|<\varepsilon$. Przez $\mathcal R[a,b]$ oznaczamy funkcje całkowalne w sensie Riemanna, a przez $\mathcal L^1$ konkretne funkcje całkowalne w sensie Lebesgue'a; $L^1$ oznacza ich klasy równości prawie wszędzie.
\subsection{Cechowania i definicja całki}
\begin{definicja}
\label{ma1:kh:def:kh-gauge}\label{it:def:kh-gauge}
Dla $a<b$ podział znaczony to skończona rodzina komórek $[x_{i-1},x_i]$, gdzie $a=x_0<x_1<\cdots<x_m=b$, ze znacznikami $t_i\in[x_{i-1},x_i]$. Jego suma ma postać
\[
 S(f,\mathcal P)=\sum_{i=1}^m f(t_i)(x_i-x_{i-1}).
\]
Podział częściowy składa się z przedziałów o dodatnich długościach, parami rozłącznych wnętrzach i znacznikach należących do odpowiednich przedziałów. Na przedziale jednopunktowym przyjmujemy całkę zero.


\emph{Cechowaniem} na $[a,b]$ nazywamy dowolną funkcję $\delta:[a,b]\to(0,\infty)$. Podział znaczony $\mathcal P=\{(t_i,[x_{i-1},x_i])\}$ jest $\delta$-drobny, jeśli $t_i\in[x_{i-1},x_i]\subset(t_i-\delta(t_i),t_i+\delta(t_i))$ dla każdego $i$.
\end{definicja}
\begin{lemat}[Cousina]\label{ma1:kh:lem:kh-cousin}\label{it:lem:kh-cousin}
Niech $a<b$. Dla każdego cechowania $\delta$ na $[a,b]$ istnieje $\delta$-drobny podział znaczony tego przedziału.
\end{lemat}
\begin{proof}
Przypuśćmy, że przedział nie ma takiego podziału. Dzielimy go na połowy; co najmniej jedna również nie ma drobnego podziału, bo dwa drobne podziały połówek można połączyć. Powtarzając, dostajemy zstępujący ciąg domkniętych przedziałów $J_n$ długości dążących do zera, z których żaden nie ma podziału drobnego. Ich przecięcie jest punktem $t$. Dla dostatecznie dużego $n$ mamy $J_n\subset(t-\delta(t),t+\delta(t))$, więc pojedynczy znaczony przedział $(t,J_n)$ jest drobny. Sprzeczność.
\end{proof}
\begin{definicja}
\label{ma1:kh:def:kh-integral}\label{it:def:kh-integral}
Funkcja $f:[a,b]\to\R$ jest całkowalna w sensie Kurzweila--Henstocka z całką $I$, jeżeli dla każdego $\varepsilon>0$ istnieje cechowanie $\delta$ takie, że dla każdego $\delta$-drobnego podziału znaczonego $\mathcal P$ zachodzi $|S(f,\mathcal P)-I|<\varepsilon$.
\end{definicja}
\begin{stwierdzenie}[Jednoznaczność]
\label{ma1:kh:prop:kh-unique}\label{it:prop:kh-unique}
Wartość całki HK jest wyznaczona jednoznacznie.
\end{stwierdzenie}
\begin{proof}
Jeśli $I,J$ spełniają definicję, dla $\varepsilon>0$ bierzemy minimum cechowań odpowiadających błędom $\varepsilon/2$ dla obu wartości. Lemat Cousina zapewnia istnienie podziału drobnego dla tego minimum. Dla jego sumy $S$ mamy $|I-J|\le |I-S|+|S-J|<\varepsilon$. Dowolność $\varepsilon$ daje $I=J$.
\end{proof}
\begin{twierdzenie}[Kryterium Cauchy'ego]\label{ma1:kh:thm:kh-cauchy}\label{it:thm:kh-cauchy}
Funkcja $f$ jest całkowalna w sensie HK wtedy i tylko wtedy, gdy dla każdego $\varepsilon>0$ istnieje cechowanie $\delta$ takie, że $|S(f,\mathcal P)-S(f,\mathcal Q)|<\varepsilon$ dla wszystkich $\delta$-drobnych podziałów $\mathcal P,\mathcal Q$.
\end{twierdzenie}
\begin{proof}
Konieczność wynika z nierówności trójkąta. Dla dostateczności wybieramy cechowania $\delta_n$ odpowiadające błędom $2^{-n}$, zastępujemy je kolejno minimami z poprzednimi i wybieramy $\delta_n$-drobne podziały $\mathcal P_n$ z lematu Cousina. Dla $m\ge n$ mamy $\delta_m\le\delta_n$, więc oba podziały $\mathcal P_n,\mathcal P_m$ są $\delta_n$-drobne i $|S(f,\mathcal P_n)-S(f,\mathcal P_m)|<2^{-n}$. Sumy tworzą zatem ciąg Cauchy'ego. Niech $I$ będzie jego granicą. Dla danego $n$ i dowolnego $\delta_n$-drobnego podziału $\mathcal P$ porównujemy jego sumę z sumami $S(f,\mathcal P_m)$, $m\ge n$, i przechodzimy z $m\to\infty$. Otrzymujemy $|S(f,\mathcal P)-I|\le2^{-n}$. Dla dowolnego $\varepsilon>0$ wybieramy $n$ z $2^{-n}<\varepsilon$; cechowanie $\delta_n$ spełnia definicję całki.
\end{proof}
\begin{twierdzenie}[Liniowość, porządek i lokalność]\label{ma1:kh:thm:kh-locality}\label{it:thm:kh-locality}
Klasa funkcji HK-całkowalnych jest przestrzenią wektorową, całka jest liniowa i monotoniczna. Jeśli $a<c<b$, to $f$ jest całkowalna w sensie HK na $[a,b]$ wtedy i tylko wtedy, gdy jest całkowalna na obu $[a,c]$ i $[c,b]$; wówczas $\int_a^b f=\int_a^c f+\int_c^b f$.
\end{twierdzenie}
\begin{proof}
Dla funkcji całkowalnych $f,g$ i liczb $\alpha,\beta\in\R$ wybieramy minimum cechowań dla $f,g$ z błędem $\varepsilon/(1+|\alpha|+|\beta|)$. Liniowość sum daje
\[
 |S(\alpha f+\beta g,\mathcal P)-(\alpha I_f+\beta I_g)|
 <\varepsilon.
\]
Jeżeli $f\le g$, to dla każdego podziału $S(f,\mathcal P)\le S(g,\mathcal P)$. Wybierając wspólne cechowanie dla błędu $\varepsilon/2$ i podział z lematu Cousina, uzyskujemy $I_f\le I_g+\varepsilon$. Dowolność $\varepsilon$ daje monotoniczność.

Jeżeli $f$ jest całkowalna na $[a,b]$, ustalamy cechowanie $\delta$ z kryterium Cauchy'ego dla błędu $\varepsilon$. Na $[a,c]$ wybieramy jeden $\delta$-drobny podział $\mathcal R$ z lematu Cousina. Dla dwóch dowolnych drobnych podziałów $\mathcal P,\mathcal Q$ przedziału $[c,b]$ ich połączenia z $\mathcal R$ są drobnymi podziałami $[a,b]$. Po odjęciu wspólnego składnika $S(f,\mathcal R)$ otrzymujemy $|S(f,\mathcal P)-S(f,\mathcal Q)|<\varepsilon$. Kryterium Cauchy'ego daje całkowalność na $[c,b]$; zamiana ról obu przedziałów daje całkowalność na $[a,c]$. Stosując ten argument dwukrotnie, otrzymujemy całkowalność na każdym podprzedziale. W drugą stronę wybieramy cechowania na obu podprzedziałach dla błędu $\varepsilon/2$, dla $t\ne c$ dodatkowo mniejsze od $|t-c|$, a w punkcie $c$ bierzemy minimum cechowań. Każdy przedział drobnego podziału przecinający $c$ ma znacznik $c$; rozcinamy go w $c$. Powstałe części dają drobne podziały z tą samą sumą, ponieważ przy rozcięciu zachowujemy znacznik $c$, a długości dodają się. Każda z dwóch sum różni się od odpowiedniej całki lokalnej o mniej niż $\varepsilon/2$. Nierówność trójkąta daje błąd mniejszy od $\varepsilon$ względem sumy całek lokalnych, co dowodzi całkowalności i addytywności. Jednoznaczność zapewnia tę samą równość także wtedy, gdy zaczynaliśmy od całkowalności na całym przedziale.
\end{proof}
\begin{stwierdzenie}
$\mathcal R[a,b]\subset\HK[a,b]$ i obie całki są równe.
\end{stwierdzenie}
\begin{proof}
Z twierdzenia \ref{ma1:thm:riemann-darboux-equivalence} dla każdego $\varepsilon>0$ istnieje $\delta_0>0$ takie, że wszystkie sumy Riemanna dla podziałów o średnicy mniejszej od $\delta_0$ różnią się od całki Riemanna o mniej niż $\varepsilon$. Wybieramy cechowanie stałe $\delta(t)=\delta_0/2$. Każda komórka drobnego podziału ma długość mniejszą od $2\delta(t)=\delta_0$, więc jego suma spełnia żądane oszacowanie.
\end{proof}
\begin{przyklad}[Funkcja Dirichleta jest HK-całkowalna]\label{ma1:kh:ex:dirichlet}
Niech $\Q\cap[0,1]=\{q_1,q_2,\dots\}$ i $D=\chi_{\Q\cap[0,1]}$. Dla $\varepsilon>0$ kładziemy $\delta(q_k)=\varepsilon 2^{-k-3}$, a w punktach niewymiernych np. $\delta(t)=1$. Grupujemy wszystkie komórki o tym samym znaczniku $q_k$. Ich wnętrza są rozłączne, a wszystkie leżą w $(q_k-\delta(q_k),q_k+\delta(q_k))$, więc ich \emph{łączna} długość nie przekracza $2\delta(q_k)$, również gdy znacznik jest wspólnym końcem dwóch komórek. Stąd dla każdego drobnego podziału $0\le S(D,\mathcal P)\le\sum_k 2\delta(q_k)<\varepsilon$. Całka HK jest równa $0$.
\end{przyklad}
\begin{lemat}[Zmiany na zbiorze miary zero]\label{ma1:kh:lem:kh-null-modification}\label{it:lem:kh-null-modification}
Jeśli $h:[a,b]\to\R$ ma skończone wartości i jest równa zeru poza zbiorem $N$ miary Lebesgue'a zero, to jej całka HK istnieje i wynosi $0$. To samo dotyczy całki McShane'a. Zmiana funkcji całkowalnej na zbiorze miary zero nie zmienia wartości żadnej z tych całek.
\end{lemat}
\begin{proof}
Rozbijamy $N$ na zbiory $N_k=\{t\in N:k-1<|h(t)|\le k\}$, $k\ge1$. Dla ustalonego $\varepsilon>0$ wybieramy otwarte pokrycia $U_k\supset N_k$ z $\lambda_1(U_k)<\varepsilon2^{-k-1}/k$. Jeśli $t\in N_k$, dobieramy $\delta(t)>0$ tak, aby $(t-\delta(t),t+\delta(t))\subset U_k$; w pozostałych punktach, również w punktach $N$ o $h(t)=0$, cechowanie jest dowolne, gdyż ich wkład do sumy wynosi zero. W podziale drobnym, również wolnym w sensie McShane'a, komórki z oznaczonymi punktami należącymi do $N_k$ mają rozłączne wnętrza i leżą w $U_k$. Ich wkład do modułu sumy jest co najwyżej $k\lambda_1(U_k)$. Sumując po $k$, dostajemy błąd mniejszy od $\varepsilon$. Jeśli dwie funkcje o skończonych wartościach różnią się tylko na zbiorze zerowym, ich różnica spełnia powyższy warunek. Liniowość całki HK (twierdzenie \ref{it:thm:kh-locality}) i całki McShane'a (twierdzenie \ref{it:thm:mcshane-basic}) daje ostatnią tezę, w obu kierunkach.
\end{proof}

\begin{center}
\begin{tikzpicture}[x=1.25cm,y=1cm,>=Stealth]
 \draw[->] (-.25,0)--(7.4,0);
 \draw[very thick,blue!65!black] (0,0)--(2,0);
 \draw[very thick,teal!70!black] (2,0)--(4.8,0);
 \draw[very thick,blue!65!black] (4.8,0)--(7,0);
 \foreach \x/\name in {0/{a=x_0},2/{x_1},4.8/{x_2},7/{x_3=b}}
   {\draw (\x,.10)--(\x,-.1) node[below] {$\name$};}
 \foreach \x/\name in {1/t_1,3.2/t_2,6/t_3}
   {\fill (\x,0) circle (2pt) node[above] {$\name$};}
 \draw[blue!65!black,<->] (-.2,.65)--(2.2,.65);
 
 \draw[teal!70!black,<->] (1.5,1.2)--(4.9,1.2);
 \draw[blue!65!black,<->] (4.6,.65)--(7.4,.65);
 \node[above] at (3.2,1.25) {$t_2\pm\delta(t_2)$};
\end{tikzpicture}

\small Rysunek: komórki podziału i otoczenia cechowania; ich rozmiary zależą od znacznika. Każda komórka musi zawierać się w swoim otwartym otoczeniu.
\end{center}
\subsection{Pochodne i lemat Saksa--Henstocka}
\begin{przyklad}[Wszystkie pochodne są całkowalne HK]
Jeśli $F:[a,b]\to\R$ jest ciągła i różniczkowalna na $(a,b)$, to funkcja $f$ równa $F'$ we wnętrzu i mająca dowolne skończone wartości w $a,b$ jest HK-całkowalna oraz $\int_a^b f(x)\,\mathrm dx=F(b)-F(a)$. Dla $\varepsilon>0$ przy każdym $t\in(a,b)$ wybieramy $\delta(t)$ tak, aby $|F(y)-F(t)-f(t)(y-t)|\le\varepsilon|y-t|/(4(b-a))$ dla bliskich $y$. Na komórce $[u,v]$ ze znacznikiem $t$ odejmujemy przyrosty $F(v)-F(t)$ oraz $F(u)-F(t)$; nierówność trójkąta daje
\[
 |F(v)-F(u)-f(t)(v-u)|
 \le \frac{\varepsilon}{4(b-a)}(|v-t|+|u-t|)
 =\frac{\varepsilon}{4(b-a)}(v-u).
\]
Dla $t=a$ i $t=b$ korzystamy tylko z ciągłości $F$: dobieramy $\delta(t)>0$ tak, aby $|F(y)-F(t)|+|f(t)||y-t|<\varepsilon/4$ przy $|y-t|<\delta(t)$, $y\in[a,b]$. Komórka ze znacznikiem $a$ lub $b$ musi mieć ten punkt jako koniec; każdego z tych znaczników można użyć co najwyżej raz. Łączny błąd komórek ze znacznikami wewnętrznymi jest nie większy niż $\varepsilon/4$, a komórek ze znacznikami brzegowymi mniejszy niż $\varepsilon/2$. Suma przyrostów $F$ teleskopuje się do $F(b)-F(a)$, więc całkowity błąd jest mniejszy od $\varepsilon$. Nie potrzeba istnienia pochodnych jednostronnych na końcach.
\end{przyklad}
\begin{przyklad}[Pochodna nieabsolutnie całkowalna]\label{ma1:kh:ex:conditional-derivative}
Niech $F(x)=x^2\cos(\pi/x^2)$ dla $x\in(0,1]$ i $F(0)=0$. Wtedy $F'(0)=0$ oraz $F'(x)=2x\cos(\pi/x^2)+(2\pi/x)\sin(\pi/x^2)$. Pochodna jest całkowalna HK i $\int_0^1F'(x)\,\mathrm dx=F(1)-F(0)=-1$. Natomiast rozszerzona całka Lebesgue'a $\int_0^1|F'(x)|\,\mathrm d\lambda_1(x)=\infty$:  składnik $2x\cos(\pi/x^2)$ ma całkowalny moduł, ograniczony przez $2x$. Dla drugiego składnika podstawienie $u=\pi/x^2$ daje
\[
 \int_c^1\frac{2\pi}{x}|\sin(\pi/x^2)|\,\mathrm dx
 =\pi\int_\pi^{\pi/c^2}\frac{|\sin u|}{u}\,\mathrm du.
\]
Na każdym przedziale $[k\pi,(k+1)\pi]$, $k\ge1$, całka po prawej bez czynnika $\pi$ jest nie mniejsza niż $2/((k+1)\pi)$. Szereg harmoniczny jest rozbieżny. Z nierówności $|A+B|\ge|B|-|A|$ wynika rozbieżność całki modułu całej pochodnej.
\end{przyklad}
\begin{lemat}[Saksa--Henstocka: uzupełnianie podziału częściowego]\label{ma1:kh:lem:kh-saks-henstock}\label{it:lem:kh-saks-henstock}\label{ma1:kh:lem:hk-uzupelnianie}\label{lem:hk-uzupelnianie}
Niech $f\in\HK[a,b]$, $F(x)=\int_a^x f$, a $\delta$ będzie cechowaniem, dla którego wszystkie $\delta$-drobne sumy na $[a,b]$ różnią się od $F(b)-F(a)$ o mniej niż $\varepsilon$. Jeśli $\mathcal Q=\{(t_j,[u_j,v_j])\}_{j=1}^m$ jest dowolnym $\delta$-drobnym podziałem częściowym z rozłącznymi wnętrzami, to
\[
\left|\sum_{j=1}^m\bigl(f(t_j)(v_j-u_j)-F(v_j)+F(u_j)\bigr)\right|\le\varepsilon.
\]
Ponadto suma modułów wyrazów w nawiasach jest nie większa niż $2\varepsilon$.
\end{lemat}
\begin{proof}
Uporządkujmy przedziały częściowego podziału. Domknięcia luk pomiędzy tymi przedziałami oraz skrajnych luk przy $a,b$ tworzą skończoną rodzinę co najwyżej $m+1$ przedziałów $J_k$ o dodatniej długości i rozłącznych wnętrzach. Wraz z komórkami częściowego podziału pokrywają $[a,b]$; luki jednopunktowe pomijamy. Dla dowolnego $\eta>0$ lokalność całki HK i lemat Cousina pozwalają wybrać na każdym $J_k$ podział znaczony, drobny względem $\delta$ oraz względem cechowania dającego błąd $\eta/(m+1)$ dla lokalnej całki po $J_k$. Połączenie tych podziałów z $\mathcal Q$ daje pełny podział $\delta$-drobny. Przyrosty $F$ na wszystkich komórkach sumują się do $F(b)-F(a)$. Błąd na $J_k$ jest mniejszy od $\eta/(m+1)$, zatem suma błędów na $\mathcal Q$ ma moduł nie większy niż $\varepsilon+\eta$. Przechodzimy z $\eta\downarrow0$.

Ten sam argument działa dla każdego podzbioru komórek $\mathcal Q$. Stosując go osobno do komórek o błędzie dodatnim i ujemnym, widzimy, że suma błędów dodatnich oraz modułów błędów ujemnych nie przekracza osobno $\varepsilon$. Ich suma jest więc nie większa niż $2\varepsilon$.
\end{proof}
\begin{lemat}[Pięciokrotne powiększenie przedziałów]\label{ma1:kh:lem:fivefold-cover}
Niech $\mathcal I$ będzie rodziną domkniętych przedziałów o dodatnich długościach ograniczonych z góry przez $L<\infty$. Istnieje co najwyżej przeliczalna podrodzina $\mathcal J$ przedziałów o parami rozłącznych wnętrzach taka, że
\[
 \bigcup_{I\in\mathcal I} I\subseteq\bigcup_{J\in\mathcal J}5J,
\]
gdzie $5J$ jest domkniętym przedziałem o tym samym środku co $J$ i pięciokrotnie większej długości.
\end{lemat}
\begin{proof}
Dzielimy rodzinę na klasy $2^{-k-1}L<|I|\le 2^{-k}L$, $k=0,1,\ldots$. W każdej kolejnej klasie wybieramy maksymalną względem inkluzji podrodzinę o rozłącznych wnętrzach, rozłącznych także z wnętrzami wcześniej wybranych przedziałów. Taka rodzina istnieje na mocy lematu Zorna, ponieważ suma łańcucha rodzin dopuszczalnych jest dopuszczalna. Suma wszystkich wybranych rodzin jest przeliczalna lub skończona: każdemu wnętrzu przypisujemy pierwszą należącą do niego liczbę wymierną w ustalonej enumeracji $\Q$, a rozłączność daje różnowartościowość tego przyporządkowania.
Jeśli $I$ z klasy $k$ nie został wybrany, jego wnętrze przecina wnętrze pewnego wybranego $J$ z klasy o indeksie co najwyżej $k$ (w przeciwnym razie można byłoby go dołączyć na etapie $k$). Wtedy $|J|>|I|/2$. Dla środków $c_I,c_J$ mamy $|c_I-c_J|<(|I|+|J|)/2$. Dla $x\in I$ stąd
\[
 |x-c_J|\le |x-c_I|+|c_I-c_J|<|I|+|J|/2<5|J|/2.
\]
Zatem $I\subset 5J$. Przedział wybrany zawiera się oczywiście we własnym pięciokrotnym powiększeniu. To dowodzi tezy także dla rodziny pustej.
\end{proof}

\begin{twierdzenie}[Prawie wszędzie pochodna pierwotnej]\label{ma1:kh:thm:kh-primitive-derivative}\label{it:thm:kh-primitive-derivative}
Jeśli $f\in\HK[a,b]$, a $F(x)=\int_a^x f$, to $F$ jest ciągła, $F'(x)=f(x)$ prawie wszędzie oraz $f$ jest mierzalna w sensie Lebesgue'a.
\end{twierdzenie}
\begin{proof}
Ustalmy $t\in[a,b]$ i $\eta>0$. Wybieramy cechowanie z lematu Saksa--Henstocka dla błędu $\eta/2$ i bierzemy $x$ z $|x-t|<\delta(t)$. Pojedyncza komórka o końcach $x,t$ ze znacznikiem $t$ jest drobnym podziałem częściowym. Oszacowanie modułu sumy błędów przez $\eta/2$ z lematu \ref{it:lem:kh-saks-henstock} daje
$|F(x)-F(t)-f(t)(x-t)|\le\eta/2$, niezależnie od kolejności $x,t$. Jeśli ponadto $|x-t|<\eta/(2(1+|f(t)|))$, to $|F(x)-F(t)|<\eta$. Dowodzi to ciągłości, także jednostronnej na końcach. Z lematu Saksa--Henstocka otrzymujemy następujące oszacowanie: dla ustalonego $r>0$ zbiór punktów, w których istnieją dowolnie małe przedziały $[u,v]$ zawierające $x$ i spełniające $|F(v)-F(u)-f(x)(v-u)|>r(v-u)$, ma miarę zewnętrzną zero. Istotnie, ustalmy dowolne $\varepsilon>0$ i wybierzmy cechowanie $\delta$ zapewniające błąd mniejszy od $\varepsilon$ w definicji całki na $[a,b]$. Lemat \ref{it:lem:kh-saks-henstock} ogranicza wówczas sumę modułów błędów każdego drobnego podziału częściowego przez $2\varepsilon$. Wybieramy tylko takie przedziały ze znacznikiem $x$, które leżą w $(x-\delta(x),x+\delta(x))$. Nadal istnieją dowolnie krótkie przedziały tego rodzaju przy każdym punkcie rozpatrywanego zbioru, więc tworzą drobne pokrycie Vitaliego. Stosujemy lemat o pięciokrotnym powiększeniu (twierdzenie \atlasref{ME-24::thm:5r-24}{ME-24: lemat pokryciowy 5r}) do wnętrz tych przedziałów. Otrzymujemy co najwyżej przeliczalną rodzinę przedziałów o rozłącznych wnętrzach. Ich pięciokrotne powiększenia pokrywają także końce wszystkich przedziałów z pierwotnej rodziny: w dowodzie lematu, jeżeli przedział o środku $z$ i promieniu $r_0$ przecina wybrany przedział o środku $w$ i promieniu $s_0>r_0/2$, to dla każdego $y$ z domknięcia pierwszego przedziału mamy $|y-w|\le r_0+|z-w|<2r_0+s_0<5s_0$. Dla przedziału wybranego jego domknięcie również mieści się w pięciokrotnym powiększeniu. W szczególności pokryty jest cały rozpatrywany zbiór punktów, nawet gdy znacznik leży na końcu komórki. Każdemu wybranemu przedziałowi przypisujemy jeden znacznik, dla którego zachodzi nierówność definiująca zbiór wyjątkowy. Dla każdej skończonej podrodziny lemat Saksa--Henstocka daje $r\sum_j(v_j-u_j)<2\varepsilon$. Biorąc supremum po skończonych podrodzinach, otrzymujemy $\sum_j(v_j-u_j)\le2\varepsilon/r$. Z przeliczalnej podaddytywności miary zewnętrznej wynika, że miara zewnętrzna zbioru jest nie większa niż $10\varepsilon/r$, a $\varepsilon$ jest dowolne. Stosujemy ten argument do $r=1/n$, $n\ge1$, i wykluczamy przeliczalną sumę otrzymanych zbiorów zerowych oraz końce $a,b$. Dla każdego pozostałego $x$ i każdego $n$ istnieje promień, poniżej którego żaden przedział zawierający $x$ nie narusza oszacowania z $r=1/n$. W szczególności przedział o końcach $x,x+h$ daje
\[
 \left|\frac{F(x+h)-F(x)}h-f(x)\right|\le1/n
\]
dla dostatecznie małych niezerowych $h$ obu znaków. To dowodzi $F'=f$ prawie wszędzie. Przedłużamy $F$ przez stałe poza $[a,b]$. Ciągłe funkcje $n(F(x+1/n)-F(x))$ zbiegają do $f(x)$ prawie wszędzie; ich granica górna jest mierzalna, a po zastąpieniu nieskończonych wartości zerem otrzymujemy skończoną funkcję mierzalną równą $f$ prawie wszędzie. Zupełność miary Lebesgue'a daje mierzalność $f$.
\end{proof}
\begin{uwaga}[Dlaczego ograniczamy każdą sumę częściową]
Nie wystarcza samo oszacowanie modułu sumy wszystkich błędów: dodatnie i ujemne błędy mogą się znosić. Dopiero możliwość wybrania dowolnego podzbioru komórek w lemacie \ref{ma1:kh:lem:hk-uzupelnianie} pozwala otrzymać sumę modułów. Jest ona potrzebna w dowodzie całkowania przez części i przy własnościach funkcji pierwotnych.
\end{uwaga}
\subsection{Całki niewłaściwe i osobliwości}
\noindent\emph{Porównanie w Bibliotece matematycznej:} \atlasref{AM1-36::ma1:kh:thm:kh-hake}{AM1-36: twierdzenie Hake'a}.

\begin{twierdzenie}[Hake'a: związek z całkami niewłaściwymi]\label{ma1:kh:thm:kh-hake}\label{it:thm:kh-hake}
Niech $f:[a,b]\to\R$ ma skończone wartości. Wówczas $f\in\HK[a,b]$ wtedy i tylko wtedy, gdy $f\in\HK[c,b]$ dla każdego $c\in(a,b)$ oraz istnieje skończona granica $L=\lim_{c\downarrow a}\int_c^b f$. W tym przypadku $\int_a^b f=L$, niezależnie od wartości $f(a)$.
\end{twierdzenie}
\begin{proof}
Jeśli $f\in\HK[a,b]$, lokalność daje całkowalność na $[c,b]$ i $\int_c^b f=F(b)-F(c)$, gdzie $F(x)=\int_a^x f$. Ciągłość $F$ w $a$ i $F(a)=0$ dają granicę $F(b)$. Dla dostateczności definiujemy $F(a)=0$ i $F(x)=L-\int_x^b f$ dla $x>a$. Funkcja jest ciągła, a jej przyrosty na przedziałach oddalonych od $a$ są lokalnymi całkami. Wybieramy ciąg $b=c_0>c_1>\dots\downarrow a$. Na $J_n=[c_n,c_{n-1}]$ wybieramy cechowanie $\delta_n$ zapewniające błąd definicyjny mniejszy od $\varepsilon2^{-n-3}$. Lemat Saksa--Henstocka daje wówczas sumę modułów błędów na podziale częściowym nie większą niż $\varepsilon2^{-n-2}$. Dla $t\in(c_n,c_{n-1})$ bierzemy minimum $\delta_n(t)$ oraz połowy odległości od obu końców. W punkcie $c_n$, $n\ge1$, bierzemy minimum $\delta_n(c_n)$, $\delta_{n+1}(c_n)$ i połowy odległości od $c_{n-1},c_{n+1}$. W punkcie $b$ wystarczy minimum $\delta_1(b)$ i $(b-c_1)/2$. Komórka przecinająca $c_n$ musi mieć znacznik $c_n$ i nie może sięgać kolejnego punktu podziału; po rozcięciu obie części są drobne względem lokalnych cechowań. W $a$ wybieramy $\delta(a)>0$ z $|F(y)|+|f(a)|(y-a)<\varepsilon/2$ dla $a<y<a+\delta(a)$, a poza $a$ zmniejszamy cechowania tak, aby przedziały z innymi znacznikami nie zawierały $a$. Przedział początkowy dowolnego drobnego podziału ma znacznik $a$, pozostałe można pogrupować według $J_n$ (i ewentualnie rozciąć w $c_n$). Rozcięcie nie zmienia sumy Riemanna ani sumy przyrostów $F$. Na każdym $J_n$ otrzymujemy podział częściowy, więc suma modułów błędów nie przekracza $\varepsilon2^{-n-2}$. Komórka początkowa wnosi błąd mniejszy od $\varepsilon/2$. Ponieważ $\sum_{n\ge1}2^{-n-2}=1/4$, całkowity błąd jest mniejszy od $\varepsilon$. Suma przyrostów wynosi $F(b)-F(a)=L$.
\end{proof}
\begin{przyklad}[Szeregi i całkowalność]
Dla $I_n=(2^{-n},2^{-n+1}]$ i ciągu $(a_n)_{n\ge1}$ kładziemy $f(x)=2^n a_n$ na $I_n$, a $f(0)=0$. Wówczas
\[
f\in\HK[0,1]\iff\sum_{n\ge1}a_n\ \text{jest zbieżny},\qquad
f\in\mathcal L^1[0,1]\iff\sum_{n\ge1}|a_n|<\infty.
\]
Na każdym $[c,1]$, $c>0$, funkcja jest schodkowa z co najwyżej skończoną liczbą skoków. Dla $c\in I_n$ całka po $[c,1]$ jest równa $\sum_{k=1}^{n-1}a_k+2^na_n(2^{-n+1}-c)$, przy czym ostatni składnik ma moduł co najwyżej $|a_n|$. Gdy szereg jest zbieżny, $a_n\to0$, więc ta całka ma granicę równą sumie szeregu. Odwrotnie, biorąc $c=2^{-n}$ otrzymujemy sumy częściowe szeregu. Twierdzenie Hake'a daje pierwszą równoważność; druga z $\int_{I_n}|f|\,\mathrm d\lambda_1=|a_n|$. Ponadto $f$ jest całkowalna Riemanna dokładnie wtedy, gdy $\sup_n2^n|a_n|<\infty$: konieczność wynika z ograniczoności funkcji całkowalnej Riemanna, a przy skończonym supremum funkcja jest ograniczona i jej punkty nieciągłości zawierają się w przeliczalnym zbiorze $\{0\}\cup\{2^{-n}:n\ge0\}$. Kryterium Lebesgue'a całkowalności Riemanna daje dostateczność. Dla $a_n=(-1)^{n+1}/n$ otrzymujemy $\int_0^1f=\log2$ w sensie HK, lecz $\int_0^1|f|\,\mathrm d\lambda_1=\infty$.
\end{przyklad}
\begin{przyklad}[Nie każda osobliwość jest całkowalna]
Funkcja $f(x)=1/x$ na $(0,1]$, niezależnie od skończonej wartości w $0$, nie jest całkowalna HK na $[0,1]$, bo $\int_c^1\mathrm dx/x=-\log c\to\infty$. Na $[-1,1]$ funkcja $1/x$ z dowolną skończoną wartością w $0$ ma wartość główną Cauchy'ego równą $0$, ale nie ma całki HK, gdyż jej ograniczenie do $[0,1]$ nie jest HK-całkowalne.
\end{przyklad}
\begin{twierdzenie}[Rodzina całek oscylacyjnych]\label{ma1:kh:tw:hk-oscylacyjne}\label{tw:hk-oscylacyjne}
Niech $p\in\R$, $q>0$ i niech $f_{p,q}:[0,1]\to\R$ będzie określona wzorem $f_{p,q}(0)=0$ oraz $f_{p,q}(x)=x^{-p}\sin(x^{-q})$ dla $x\in(0,1]$. Wówczas
\[
 f_{p,q}\in\HK[0,1]\quad\Longleftrightarrow\quad p<q+1,
 \qquad f_{p,q}\in\mathcal L^1[0,1]\quad\Longleftrightarrow\quad p<1.
\]
W szczególności dla $1\le p<q+1$ całka HK istnieje, lecz nie jest bezwzględnie zbieżna.
\end{twierdzenie}
\begin{proof}
Na $[\varepsilon,1]$ funkcja jest ciągła. Podstawiając $t=x^{-q}$, uzyskujemy
\[
\int_\varepsilon^1x^{-p}\sin(x^{-q})\,\mathrm dx
=\frac1q\int_1^{\varepsilon^{-q}}t^{(p-1)/q-1}\sin t\,\mathrm dt.
\]
Oznaczmy wykładnik przez $\beta=(p-1)/q-1$. Gdy $\beta<0$, funkcja $t^\beta$ maleje do zera. Dla $1\le A<B$ całkowanie przez części daje
\[
\int_A^B t^\beta\sin t\,\mathrm dt
=[-t^\beta\cos t]_A^B+\beta\int_A^B t^{\beta-1}\cos t\,\mathrm dt.
\]
Moduł prawej strony jest nie większy niż $A^\beta+B^\beta+|\beta|\int_A^B t^{\beta-1}\,\mathrm dt=2A^\beta$. Zatem kryterium Cauchy'ego daje zbieżność całki niewłaściwej. Gdy $\beta=0$, pierwotna $-\cos t$ nie ma granicy. Jeśli $\beta>0$, na półokresie $[k\pi,(k+1)\pi]$, $k\ge1$, sinus ma stały znak. Stąd
\[
 \left|\int_{k\pi}^{(k+1)\pi}t^\beta\sin t\,\mathrm dt\right|
 \ge 2(k\pi)^\beta,
\]
gdyż całka z $|\sin t|$ po tym półokresie wynosi $2$. Narusza to kryterium Cauchy'ego całki niewłaściwej. Twierdzenie Hake'a daje pierwszą równoważność, ponieważ $\beta<0$ dokładnie wtedy, gdy $p<q+1$.

Dla $\beta<-1$ całka modułu jest skończona przez oszacowanie $t^\beta|\sin t|\le t^\beta$ dla $t\ge1$. Dla $-1\le\beta<0$ mamy na każdym półokresie
\[
 \int_{k\pi}^{(k+1)\pi}t^\beta|\sin t|\,\mathrm dt
 \ge 2((k+1)\pi)^\beta.
\]
Szereg tych dolnych ograniczeń jest rozbieżny, ponieważ $(k+1)^\beta\ge 1/(k+1)$. Dla $\beta\ge0$ dolne ograniczenie wynosi $2(k\pi)^\beta\ge2$ i również daje rozbieżność. Zatem bezwzględna całkowalność zachodzi dokładnie dla $\beta<-1$, czyli $p<1$.
\end{proof}
\begin{przyklad}[Dokładny próg przy całce $x^{-2}\sin(x^{-2})$]
Dla $p=q=2$ mamy $2<3$, więc całka $\int_0^1x^{-2}\sin(x^{-2})\,\mathrm dx$ istnieje w sensie HK, ale nie w sensie Lebesgue'a. Podstawienie z dowodu daje użyteczną reprezentację obliczeniową
$\int_0^1x^{-2}\sin(x^{-2})\,\mathrm dx=\tfrac12\int_1^\infty t^{-1/2}\sin t\,\mathrm dt$.
\end{przyklad}
\subsection{Podstawienie oraz całkowanie przez części}
\begin{twierdzenie}[Mnożniki o ograniczonej wariacji]\label{ma1:kh:thm:bv-multiplier}
Niech $f\in\HK[a,b]$, $F(x)=\int_a^x f$, a $g:[a,b]\to\R$ ma ograniczoną wariację. Wtedy $fg\in\HK[a,b]$ oraz
\[
\int_a^b f(x)g(x)\,\mathrm dx=F(b)g(b)-F(a)g(a)-\int_a^b F\,\mathrm dg,
\]
gdzie ostatnia całka jest zwykłą całką Riemanna--Stieltjesa.
\end{twierdzenie}
\begin{proof}
Najpierw uzasadnimy istnienie całki Riemanna--Stieltjesa. Dla podziału $P$ ze znacznikami definiujemy sumę $\sum_i F(t_i)(g(x_i)-g(x_{i-1}))$. Niech $V=\Var_{[a,b]}g$ oraz $\omega_F(s)=\sup_{|x-y|\le s}|F(x)-F(y)|$ dla $x,y\in[a,b]$. Jeśli $P'$ jest zagęszczeniem $P$, dowolna suma na $P'$ różni się od danej sumy na $P$ o co najwyżej $\omega_F(\mesh P)V$: rozpisujemy przyrosty $g$ jako sumy przyrostów na mniejszych komórkach i stosujemy nierówność trójkąta. Dla dwóch podziałów o średnicach mniejszych od $s$ porównanie ze wspólnym zagęszczeniem daje różnicę nie większą niż $2\omega_F(s)V$. Jednostajna ciągłość $F$ daje $\omega_F(s)\to0$. Wybierając dowolny ciąg podziałów o średnicach dążących do zera, otrzymujemy ciąg Cauchy'ego sum, a porównanie z jego granicą dowodzi istnienia wspólnej granicy wszystkich sum. Nazywamy ją całką Riemanna--Stieltjesa $\int_a^b F\,\mathrm dg$. Dla $\delta$ z lematu Saksa--Henstocka i $\Delta F_i=F(x_i)-F(x_{i-1})$ mamy
$|S(fg,\mathcal P)-\sum_i g(t_i)\Delta F_i|\le2\varepsilon\|g\|_\infty$.
Ponieważ $\sum_i|g(t_i)-g(x_{i-1})|\le \Var_{[a,b]}g$ i $|\Delta F_i|\le\omega_F(\mesh P)$, różnica między $\sum_i g(t_i)\Delta F_i$ i $\sum_i g(x_{i-1})\Delta F_i$ nie przekracza $\omega_F(\mesh P)\Var_{[a,b]}g$. Dokładne sumowanie przez części daje
$\sum_i g(x_{i-1})\Delta F_i=F(b)g(b)-F(a)g(a)-\sum_iF(x_i)(g(x_i)-g(x_{i-1}))$.
Ostatnia suma dąży do $\int F\,\mathrm dg$. Bierzemy $\varepsilon$ małe i dodatkowo ograniczamy cechowanie stałą zapewniającą dostatecznie małą średnicę podziałów.
\end{proof}
\begin{uwaga}[Ograniczenia całki HK]
Klasa $\HK[a,b]$ nie jest kratą: z całkowalności $f$ nie wynika całkowalność $|f|$. Gdyby moduł pochodnej z przykładu o funkcji $x^2\cos(\pi/x^2)$ był HK-całkowalny, to na mocy twierdzenia \ref{it:thm:hk-absolute-integrability} byłby całkowalny w sensie Lebesgue'a, wbrew wykazanej rozbieżności jego całki. Nie jest też algebrą, ponieważ $x^{-1/2}$ na $(0,1]$ po odpowiednim określeniu w zerze ma całkę HK, a jej kwadrat $1/x$ nie. Seminorma Alexiewicza $\|f\|_A=\sup_{x\in[a,b]}\left|\int_a^x f\right|$ jest normą po utożsamieniu funkcji równych prawie wszędzie, lecz otrzymana przestrzeń nie jest zupełna (dowód i opis uzupełnienia podano w stwierdzeniu o normie Alexiewicza w niniejszym rozdziale).
\end{uwaga}
\begin{twierdzenie}[Podstawienie przez dyfeomorfizm klasy $C^1$]\label{ma1:kh:tw:hk-podstawienie}\label{tw:hk-podstawienie}
Niech $a<b$, $u<v$, $f\in\HK[a,b]$ i niech $\phi:[u,v]\to[a,b]$ będzie rosnącym dyfeomorfizmem klasy $C^1$, to znaczy $\phi\in C^1[u,v]$, $\phi(u)=a$, $\phi(v)=b$ i $\phi'(t)>0$ dla $t\in[u,v]$. Wówczas $t\mapsto f(\phi(t))\phi'(t)$ jest HK-całkowalna na $[u,v]$ oraz
\[
\int_u^v f(\phi(t))\phi'(t)\,\mathrm dt=\int_a^b f(x)\,\mathrm dx.
\]
Dla malejącego dyfeomorfizmu z $\phi(u)=b$, $\phi(v)=a$ i $\phi'<0$ całka po prawej ma odwrócone granice.
\end{twierdzenie}
\begin{proof}
Niech $F(x)=\int_a^x f$. Funkcja $F$ jest ciągła i $F'=f$ prawie wszędzie. Najpierw uzasadnimy, że $F\circ\phi$ należy do $\AC G_*[u,v]$. Z charakteryzacji pierwotnych HK (twierdzenie \ref{it:thm:kh-acgstar-characterization} w rozdziale o wąskiej całce Denjoya) mamy $[a,b]=\bigcup_n E_n$, gdzie zbiory $E_n$ są domknięte i $F\in\AC_*(E_n)$. Zbiory $\phi^{-1}(E_n)$ są domknięte i pokrywają $[u,v]$. Ponieważ $\phi$ jest lipschitzowska, dla rozłącznych przedziałów o końcach w $\phi^{-1}(E_n)$ i o małej łącznej długości obrazy mają także małą sumę długości i rozłączne wnętrza. Równość
$\osc_J(F\circ\phi)=\osc_{\phi(J)}F$ dowodzi $F\circ\phi\in\AC_*(\phi^{-1}(E_n))$.

Dyfeomorfizm i jego odwrotność są lipschitzowskie na zwartych przedziałach, więc przeciwobraz zbioru miary zero ma miarę zero. W punktach, w których $F'(\phi(t))=f(\phi(t))$, zwykła reguła łańcuchowa daje $(F\circ\phi)'(t)=f(\phi(t))\phi'(t)$. Twierdzenie \ref{it:thm:acgstar-derivative-integrable} daje całkowalność HK tej pochodnej i wzór Newtona--Leibniza; ostatni przyrost wynosi $F(b)-F(a)$. Dla malejącego dyfeomorfizmu argument jest taki sam: obrazy komórek zapisujemy z końcami w kolejności rosnącej, ich długości kontroluje $\max|\phi'|$, a odwrotność jest lipschitzowska z oszacowania $\min|\phi'|>0$. Przyrost funkcji złożonej wynosi tym razem $F(a)-F(b)$.
\end{proof}
\begin{uwaga}[Zakres twierdzenia]
Warunek $\phi'(t)>0$ zabezpiecza przed trudnościami z przenoszeniem zbiorów wyjątkowych. Istnieją mocniejsze twierdzenia o podstawieniu dla mniej regularnych funkcji $\phi$, ale bez dodatkowych warunków na składanie pierwotnych i na zbiory wyjątkowe nie wolno automatycznie zastępować dyfeomorfizmu dowolną funkcją bezwzględnie ciągłą.
\end{uwaga}
\begin{wniosek}[Całkowanie przez części z mnożnikiem bezwzględnie ciągłym]\label{ma1:kh:wn:hk-ibp-ac}\label{wn:hk-ibp-ac}
Niech $f\in\HK[a,b]$, $F(x)=\int_a^x f$, a $g\in\AC[a,b]$. Wtedy $fg\in\HK[a,b]$ i
\[
\int_a^b f(x)g(x)\,\mathrm dx=F(b)g(b)-F(a)g(a)-\int_a^b F(x)g'(x)\,\mathrm d\lambda_1(x).
\]
Ostatnia całka jest całką Lebesgue'a, ponieważ $F$ jest ograniczona i $g'\in L^1[a,b]$.
\end{wniosek}
\begin{proof}
Funkcja $g$ ma wahanie ograniczone, gdyż $\Var(g)\le\int_a^b|g'|\,\mathrm d\lambda_1$. Stosujemy twierdzenie o mnożnikach z tej sekcji. Z twierdzenia \ref{ma1:kh:thm:ac-lebesgue-primitive} mamy $g(v)-g(u)=\int_u^v g'\,\mathrm d\lambda_1$. Różnica sumy Stieltjesa i całki Lebesgue'a jest zatem ograniczona przez
\[
 \sum_i\int_{x_{i-1}}^{x_i}|F(t_i)-F(x)|\,|g'(x)|\,\mathrm d\lambda_1(x)
 \le \omega_F(\mesh P)\int_a^b|g'|\,\mathrm d\lambda_1.
\]
Prawa strona dąży do zera, co daje żądaną identyfikację.
\end{proof}
\begin{przyklad}[Warunkowo zbieżna całka niewłaściwa]
Niech $f(x)=\sin(1/x)/x$ dla $x\in(0,1]$ i $f(0)=0$. Podstawienie $t=1/x$ na $[\varepsilon,1]$ daje
$\int_\varepsilon^1f(x)\,\mathrm dx=\int_1^{1/\varepsilon}(\sin t)/t\,\mathrm dt$.
Całka po prawej jest zbieżna na mocy kryterium Dirichleta, zatem twierdzenie Hake'a daje
$\int_0^1f(x)\,\mathrm dx=\int_1^\infty (\sin t)/t\,\mathrm dt$ w sensie HK. Nie zachodzi całkowalność Lebesgue'a wartości bezwzględnej, gdyż $\int_1^\infty|\sin t|/t\,\mathrm dt=\infty$.
\end{przyklad}
\subsection{Twierdzenia o zbieżności}
\begin{twierdzenie}[Zbieżność z majorantą Lebesgue'a]\label{ma1:kh:tw:hk-zbieznosc-dct}\label{tw:hk-zbieznosc-dct}
Niech $f_n:[a,b]\to\R$ będą mierzalne i HK-całkowalne, $f_n\to f$ prawie wszędzie dla $f:[a,b]\to\R$, a $g\in L^1[a,b]$ niech spełnia $|f_n|\le g$ prawie wszędzie dla każdego $n$. Wówczas $f,f_n\in\mathcal L^1[a,b]$ i
$\lim_{n\to\infty}\int_a^bf_n\,\mathrm dx=\int_a^bf\,\mathrm dx$, przy czym po obu stronach całkę można rozumieć zarówno w sensie Lebesgue'a, jak i HK.
\end{twierdzenie}
\begin{proof}
Wybieramy skończonego mierzalnego reprezentanta $g$ klasy $L^1$. Poza jednym zbiorem miary zero (sumą zbiorów wyjątkowych dla wszystkich nierówności i zbieżności) zachodzą jednocześnie $|f_n|\le g$ i $f_n\to f$, stąd $|f|\le g$. Funkcja $f$ jest mierzalna, ponieważ poza tym zbiorem jest granicą funkcji mierzalnych, a miara Lebesgue'a jest zupełna. Oszacowania przez $g$ dają bezwzględną całkowalność Lebesgue'a każdej $f_n$ i $f$. Stosujemy twierdzenie \ref{thm:dct} oraz zgodność całek Lebesgue'a i HK (twierdzenie \ref{it:thm:mcshane-lebesgue} i inkluzja całki McShane'a w HK). Nieujemna majoranta HK-całkowalna jest automatycznie całkowalna w sensie Lebesgue'a na mocy kryterium bezwzględnej całkowalności. Dwustronne ograniczenie przez funkcje HK-całkowalne, które same nie muszą należeć do $L^1$, oraz zbieżność monotoniczną omawiamy w niniejszym rozdziale (twierdzenie \ref{it:thm:hk-limit-theorems}).
\end{proof}
\begin{przyklad}[Dlaczego sama zbieżność punktowa nie wystarcza]
Niech $f_n(0)=0$ i $f_n(x)=n\chi_{(0,1/n]}(x)$ na $[0,1]$. Każda $f_n$ jest całkowalna w sensie Riemanna, a więc HK-całkowalna, oraz $f_n(x)\to0$ w każdym punkcie $[0,1]$. Mimo to $\int_0^1f_n\,\mathrm dx=1$ dla każdego $n$. Nie istnieje jedna całkowalna w sensie Lebesgue'a funkcja $g$ z $|f_n|\le g$ prawie wszędzie dla wszystkich $n$: w przeciwnym razie twierdzenie o zbieżności zmajoryzowanej wymuszałoby granicę całek równą zeru.
\end{przyklad}
\begin{stwierdzenie}[Jednostajna zbieżność]\label{ma1:kh:st:hk-zbieznosc-jednostajna}\label{st:hk-zbieznosc-jednostajna}
Jeśli $f_n\in\HK[a,b]$ i $f_n\to f$ jednostajnie, to $f\in\HK[a,b]$ oraz $\int_a^bf_n\to\int_a^bf$.
\end{stwierdzenie}
\begin{proof}
Dla dwóch funkcji HK-całkowalnych $h,k$ o ograniczonej różnicy zachodzi $|\int_a^b(h-k)|\le(b-a)\|h-k\|_\infty$, ponieważ nierówność dla sum przechodzi na granicę. Dla dostatecznie dużych $n,m$ różnica $f_n-f_m$ jest ograniczona i jej norma supremum dąży do zera. Stąd całki $I_n=\int_a^bf_n$ tworzą ciąg Cauchy'ego. Dla ustalonego $\varepsilon>0$ wybieramy $n$ z $(b-a)\|f-f_n\|_\infty<\varepsilon/3$ oraz $|I-I_n|<\varepsilon/3$, gdzie $I=\lim I_n$. Następnie bierzemy cechowanie dla $f_n$ zapewniające błąd sumy mniejszy niż $\varepsilon/3$. Nierówność trójkąta kończy dowód bez wcześniejszego zakładania całkowalności $f$.
\end{proof}


\begin{twierdzenie}[Jednostajna całkowalność względem cechowania]
\label{ma1:kh:thm:kh-equiintegrability}\label{it:thm:kh-equiintegrability}
Niech $f_n\in\HK[a,b]$ i $f_n(t)\to f(t)\in\R$ dla każdego $t\in[a,b]$. Załóżmy, że dla każdego $\varepsilon>0$ istnieje jedno cechowanie $\delta$ takie, że
\[
 |S(f_n,\mathcal P)-I_n|<\varepsilon
 \quad\text{dla wszystkich }n\ge1
 \text{ i wszystkich }\delta\text{-drobnych }\mathcal P,
 \qquad I_n=\int_a^b f_n.
\]
Wtedy $f\in\HK[a,b]$ oraz $I_n\to\int_a^b f$.
\end{twierdzenie}
\begin{proof}
Dla $\varepsilon>0$ wybieramy wspólne cechowanie dla błędu $\varepsilon/3$ i ustalamy jeden drobny podział $\mathcal P$ z lematu Cousina. Ma on skończenie wiele znaczników, więc $S(f_n,\mathcal P)\to S(f,\mathcal P)$ dzięki zbieżności punktowej. Dla dostatecznie dużych $n,m$ różnica tych dwóch sum ma moduł mniejszy od $\varepsilon/3$. Nierówność trójkąta daje $|I_n-I_m|<\varepsilon$. Zupełność $\R$ zapewnia granicę $I$. Teraz dla dowolnej dokładności $\eta>0$ bierzemy wspólne cechowanie dla błędu $\eta/2$. Dla każdego jego drobnego podziału $\mathcal Q$ przechodzimy z $n$ do nieskończoności w nierówności $|S(f_n,\mathcal Q)-I_n|<\eta/2$. Otrzymujemy $|S(f,\mathcal Q)-I|\le\eta/2<\eta$, zatem $f$ jest całkowalna z całką $I$. Nie zakładaliśmy z góry istnienia całki z granicy.
\end{proof}

\begin{uwaga}[Całka Henstocka--Kurzweila--Stieltjesa]
Zastępując w sumach drobnych $x_i-x_{i-1}$ przez $\alpha(x_i)-\alpha(x_{i-1})$, otrzymujemy całkę HK--Stieltjesa. W szczególności dla $f=\alpha=\ind{[c,b]}$, $a<c<b$, całka istnieje i wynosi $1$. Cechowanie można dobrać tak, aby każdy przedział o niezerowym przyroście $\alpha$ miał znacznik $c$. Reguła całkowania przez części wymaga tu poprawek na wspólne skoki.
\end{uwaga}

\subsection{Norma Alexiewicza i twierdzenia graniczne dla HK}
\begin{stwierdzenie}[Uzupełnienie w normie Alexiewicza]\label{it:thm:alexiewicz-completion}
Niech $a<b$. Przestrzeń funkcji HK-całkowalnych modulo równość prawie wszędzie, z normą $\|[f]\|_A=\sup_{x\in[a,b]}|(\HK)\int_a^x f|$, nie jest zupełna. Jej uzupełnienie izometryczne można utożsamić z $C_0[a,b]=\{F\in C[a,b]:F(a)=0\}$ z normą supremum (symbol $C_0[a,b]$ w tym stwierdzeniu oznacza funkcje zerujące się w $a$, nie przestrzeń funkcji znikających w nieskończoności).
\end{stwierdzenie}
\begin{proof}
Przyporządkowanie $f\mapsto F(x)=(\HK)\int_a^x f$ przyjmuje wartości w $C_0[a,b]$: pierwotna jest ciągła, a jej wartość w $a$ wynosi zero. Wynika to z twierdzenia \ref{it:thm:kh-primitive-derivative}. Jest liniowe dzięki liniowości i lokalności HK, a jego jądro stanowią dokładnie funkcje zerowe prawie wszędzie: jeden kierunek wynika z lematu \ref{ma1:kh:lem:kh-null-modification}, drugi z twierdzenia \ref{ma1:kh:thm:kh-primitive-derivative}. Na przestrzeni ilorazowej jest zatem różnowartościową izometrią, ponieważ $\|F\|_\infty=\|f\|_A$. W szczególności $\|\cdot\|_A$ jest na niej normą. Jego obraz zawiera wszystkie funkcje przedziałami liniowe zerujące się w $a$: dla skończonego podziału wybieramy funkcję podcałkową równą nachyleniu na wnętrzu każdej komórki, a na końcach przyjmujemy dowolne skończone wartości. Jest to ograniczona mierzalna funkcja całkowalna w sensie Lebesgue'a, więc z twierdzenia \ref{it:thm:mcshane-lebesgue} i inkluzji McShane'a w HK ma tę samą całkę w sensie HK. Dokładniej, dla węzłów $a=x_0<\cdots<x_m=b$ oraz zadanej funkcji przedziałami liniowej $P$ o $P(a)=0$ wybieramy $p_j=(P(x_j)-P(x_{j-1}))/(x_j-x_{j-1})$. Dla $x\in[x_{k-1},x_k]$ zgodność całek na podprzedziałach i zerowa miara końców dają
\[
 (\HK)\int_a^x f
 =\sum_{j<k}p_j(x_j-x_{j-1})+p_k(x-x_{k-1})
 =P(x_{k-1})-P(a)+P(x)-P(x_{k-1})=P(x).
\]
 Dla $F\in C_0[a,b]$ i $\varepsilon>0$ z jednostajnej ciągłości $F$ na zwartym przedziale wybieramy podział o tak małej średnicy, że oscylacja $F$ na każdej komórce jest mniejsza od $\varepsilon$. Interpolant liniowy $P$ wartości $F$ na końcach komórek zeruje się w $a$. Dla $x$ w komórce $[u,v]$ mamy $P(x)=(1-\theta)F(u)+\theta F(v)$, gdzie $\theta=(x-u)/(v-u)$, zatem $|P(x)-F(x)|\le (1-\theta)|F(u)-F(x)|+\theta|F(v)-F(x)|<\varepsilon$. Domknięciem obrazu jest więc całe $C_0[a,b]$. Nie każda funkcja z tej przestrzeni jest pierwotną HK: funkcja $x\mapsto C((x-a)/(b-a))$ dla funkcji Cantora $C$ z przykładu \ref{ex:cantor-measure-17} jest ciągła, zeruje się w $a$ i przyjmuje wartość $1$ w $b$. Poza afinicznym przeciwobrazem zbioru Cantora jest lokalnie stała, a więc ma pochodną zero. Przeciwobraz tego zbioru ma miarę zero na mocy \atlasref{ME-13::eq:linear-volume-13}{ME-13: wzór na zmianę objętości}. Gdyby była pierwotną HK funkcji $f$, twierdzenie \ref{it:thm:kh-primitive-derivative} dawałoby $f=0$ prawie wszędzie, więc lemat \ref{it:lem:kh-null-modification} zastosowany do ograniczenia $f$ na każdym $[a,x]$, $x>a$, wymuszałby zerową pierwotną (w $a$ również równą zero). Sprzeczność. Z twierdzenia o zupełności $C(K)$ (stwierdzenie \ref{prop:ck-banach12}) przestrzeń $C[a,b]$ jest przestrzenią Banacha z normą supremum. Ewaluacja $F\mapsto F(a)$ jest liniowa i ciągła, ponieważ $|F(a)|\le\|F\|_\infty$. Jej jądro $C_0[a,b]$ jest zatem domknięte i jako domknięta podprzestrzeń przestrzeni Banacha jest zupełne. Obraz jest gęsty, lecz nie jest całym $C_0[a,b]$, a więc nie jest domknięty i nie jest zupełny.
\end{proof}
\begin{twierdzenie}[Twierdzenia graniczne dla HK]\label{it:thm:hk-limit-theorems}\label{ma1:kh:thm:kh-order-limits}
Niech $(f_n)$ będzie ciągiem rzeczywistych funkcji
HK-całkowalnych na $[a,b]$, a $f:[a,b]\to\R$.
W przypadku granicy określonej tylko prawie wszędzie
wybieramy dowolne skończone wartości na zbiorze
wyjątkowym. Wszystkie całki bez oznaczenia miary
w poniższych trzech warunkach rozumiemy w sensie HK.
\begin{enumerate}[label=(\alph*)]
\item Jeśli $f_n\uparrow f$ prawie wszędzie i $\sup_n(\HK)\int_a^b f_n<\infty$, to $f$ jest HK-całkowalna oraz $\int_a^b f_n\to\int_a^b f$.
\item Jeśli $f_n\to f$ prawie wszędzie i istnieją $g,h\in\HK[a,b]$ takie, że $g\le f_n\le h$ prawie wszędzie dla wszystkich $n$, to $f\in\HK[a,b]$ i $\int_a^b f_n\to\int_a^b f$.
\item Jeśli $f_n\downarrow f$ prawie wszędzie i
$\inf_n(\HK)\int_a^b f_n>-\infty$, to $f\in\HK[a,b]$
i całki zbiegają do całki $f$.
\end{enumerate}
We wszystkich trzech przypadkach $f_n-f\in\mathcal L^1$
oraz dla pierwotnych $F_n(x)=(\HK)\int_a^x f_n$,
$F(x)=(\HK)\int_a^x f$ zachodzi
\[
 \|[f_n-f]\|_A
 =\|F_n-F\|_\infty
 \le V_a^b(F_n-F)
 =\int_{[a,b]}|f_n-f|\,\mathrm d\lambda_1\longrightarrow0.
\]
Nie zakładamy, że same $f_n$ lub $f$ są całkowalne
w sensie Lebesgue'a.
\end{twierdzenie}
\begin{proof}
W każdej z trzech części zbieramy wszystkie zbiory
wyjątkowe (także dla poszczególnych $n$) w jeden
mierzalny zbiór miary zero. Na tym zbiorze zastępujemy
wszystkie rozpatrywane funkcje zerem. Z lematu
\ref{ma1:kh:lem:kh-null-modification} wynika, że taka
zmiana zachowuje całkowalność HK i wartość całki
w obie strony: jeżeli dopiero zmodyfikowana funkcja
graniczna okaże się całkowalna, ten sam lemat
pozwoli przywrócić jej pierwotne wartości.
Możemy zatem zakładać punktowe nierówności
i punktową zbieżność na całym przedziale.

(a) Funkcje $u_n=f_n-f_1$ są nieujemne i HK-całkowalne, zatem należą do $\mathcal L^1$ na mocy twierdzenia \ref{it:thm:hk-absolute-integrability}. Rosną punktowo do $u=f-f_1$, więc $u$ jest mierzalna. Z liniowości HK i zgodności z całką Lebesgue'a wynika
\[
 \sup_n\int_{[a,b]}u_n\,\mathrm d\lambda_1
 =\sup_n(\HK)\int_a^b f_n-(\HK)\int_a^b f_1<\infty.
\]
Twierdzenie \ref{thm:mct} daje $u\in\mathcal L^1$ i zbieżność całek. Ponieważ $f=f_1+u$, liniowość HK daje całkowalność $f$ oraz żądaną zbieżność, po dodaniu $(\HK)\int_a^b f_1$.

(b) Funkcja $h-g$ jest nieujemna i HK-całkowalna, więc należy do $\mathcal L^1$ z twierdzenia \ref{it:thm:hk-absolute-integrability}. Także $u_n=f_n-g$ jest nieujemna, HK-całkowalna i ograniczona przez $h-g$. Jej granica punktowa $u=f-g$ jest mierzalna. Twierdzenie \ref{thm:dct} i wniosek \ref{cor:dct-l1} dają $u\in\mathcal L^1$ i $\int_{[a,b]}|u_n-u|\,\mathrm d\lambda_1\to0$. Z $f=g+u$ otrzymujemy całkowalność HK i zbieżność całek. W obu częściach przywrócenie początkowych wartości na zbiorze zerowym jest dozwolone przez lemat \ref{it:lem:kh-null-modification}. Ponadto w (a) $\int|u-u_n|\,\mathrm d\lambda_1=\int u\,\mathrm d\lambda_1-\int u_n\,\mathrm d\lambda_1\to0$, a w (b) zbieżność w $L^1$ już wykazano. W obu przypadkach $f_n-f=u_n-u$ należy do $\mathcal L^1$,
także po przywróceniu wartości na zbiorze zerowym,
gdyż miara Lebesgue'a jest zupełna.
W (c) stosujemy (a) do $-f_n\uparrow-f$; warunek
ograniczenia całek od dołu dla $f_n$ jest dokładnie
warunkiem ograniczenia ich przeciwieństw od góry.
Otrzymujemy także $\int|f_n-f|\,\mathrm d\lambda_1\to0$.

W każdym przypadku zgodność całek Lebesgue'a i HK
(twierdzenie \ref{it:thm:mcshane-lebesgue} i inkluzja
McShane'a w HK) stosujemy do różnicy na każdym
podprzedziale $[a,x]$. Zatem
$F_n(x)-F(x)=\int_{[a,x]}(f_n-f)\,\mathrm d\lambda_1$.
Jest to funkcja bezwzględnie ciągła na mocy
twierdzenia \atlasref{ME-25::thm:FTC-lebesgue-25}{ME-25: podstawowe twierdzenie rachunku całkowego}, zerująca się w $a$,
a w szczególności ma ograniczone wahanie, jak wykazano
w dowodzie twierdzenia \ref{it:thm:ac-lebesgue-primitive}.
Jej miarą Stieltjesa ze znakiem jest
$(f_n-f)\lambda_1|_{(a,b]}$: obie miary mają
te same przyrosty na przedziałach $(s,t]$.
Z jednoznaczności i równości wahania w twierdzeniu
\ref{it:thm:stieltjes-signed-bv} otrzymujemy
$V_a^b(F_n-F)=\int_{[a,b]}|f_n-f|\,\mathrm d\lambda_1$.
Istotnie wahanie miary o rzeczywistej gęstości $q$
jest miarą $|q|\lambda_1$: górne oszacowanie wynika
z nierówności trójkąta na każdym mierzalnym podziale,
a dolne z podziału na $\{q\ge0\}$ i $\{q<0\}$.
Ponieważ $F_n(a)-F(a)=0$, norma supremum różnicy
nie przekracza jej wahania. To daje wszystkie
deklarowane równości i oszacowania.
\end{proof}
\begin{przyklad}[Konieczność założeń]
Dla $f_n=n\chi_{(0,1/n)}$ mamy $f_n\to0$ punktowo na $[0,1]$, lecz $\int_0^1f_n=1$ dla wszystkich $n$; brak wspólnej całkowalnej majoranty. Z kolei $g_n(x)=\min\{n,1/x\}$ dla $x>0$, $g_n(0)=0$, rośnie punktowo do funkcji równej $1/x$ na $(0,1]$, a $0$ w zerze, i $\int_0^1g_n=1+\log n\to\infty$. Sama monotoniczność bez skończonej granicy całek nie wystarcza.
\end{przyklad}

